\exercisesection{\S~5}

\begin{enumerate}
\def\labelenumi{\arabic{enumi})}
\item
  Show that the solvable Lie algebra considered in Chap. I, §5, Exerc. 6 is not isomorphic to any decomposable Lie algebra.
\item
  Let u (resp. v) be a non-zero semi-simple (resp. nilpotent) endomorphism of V. Then the map \(\lambda u \mapsto \lambda v (\lambda \in k)\) is an isomorphism from g = ku to \(g' = kv\) which does not take semi-simple elements to semi-simple elements.
\item
  Let \(u\) be an endomorphism of \(V\) that is neither semi-simple nor nilpotent. Then \(\mathfrak{g} = ku\) is non-decomposable, but \(\operatorname{ad}_{\mathfrak{g}}\mathfrak{g}\) is decomposable.
\end{enumerate}

¶ 4) Let g be a decomposable Lie subalgebra of gl(V). Let q be a Lie subalgebra of g whose identity representation is semi-simple. There exists \(a \in \operatorname{Aut}_{e}(\mathfrak{g})\) such that \(a(\mathfrak{q})\) is contained in the subalgebra m of Prop. 7. (Imitate the proof of Prop. 7 (ii).)

¶ 5) Let g be a Lie subalgebra of gl(V). Then g is said to be algebraic if, for all \(x \in \mathfrak{g}\) , the replicas of \(x\) (Chap. I, §5, Exerc. 14) belong to g. Such an algebra is decomposable.

\begin{enumerate}
\def\labelenumi{\alph{enumi})}
\tightlist
\item
  Denote by \(a(\mathfrak{g})\) the smallest algebraic subalgebra of \(\mathfrak{gl}(\mathrm{V})\) containing \(\mathfrak{g}\) . Then
\end{enumerate}

\[
a (\mathfrak {g}) \supset e (\mathfrak {g}) \supset \mathfrak {g}.
\]

Give an example where \(a(\mathfrak{g})\) and \(e(\mathfrak{g})\) are distinct (take V of dimension 2 and \(\mathfrak{g}\) of dimension 1).

\begin{enumerate}
\def\labelenumi{\alph{enumi})}
\setcounter{enumi}{1}
\item
  Show that, if \(\mathfrak{n}\) is an ideal of \(\mathfrak{g}\) , \(\mathfrak{n}\) and \(a(\mathfrak{n})\) are ideals of \(a(\mathfrak{g})\) , and \([a(\mathfrak{g}), a(\mathfrak{n})] = [\mathfrak{g}, \mathfrak{n}]\) (imitate the proof of Prop. 4). Deduce that \(\mathcal{D}^i a(\mathfrak{g}) = \mathcal{D}^i \mathfrak{g}\) for \(i \geq 1\) and that \(\mathcal{C}^i a(\mathfrak{g}) = \mathcal{C}^i \mathfrak{g}\) for \(i \geq 2\) .
\item
  Show that every Lie algebra consisting of nilpotent elements is algebraic. \(^{8}\)
\end{enumerate}

¶ 6) Let g be a semi-simple subalgebra of \(\mathfrak{gl}(V)\) , \(\mathbf{T}(V) = \bigoplus_{n=0}^{\infty} \mathbf{T}^{n}(V)\) the tensor algebra of V, \(\mathbf{T}(V)^{\mathfrak{g}}\) the set of elements of \(\mathbf{T}(V)\) invariant under g (cf.~Chap. III, App.) and \(\overline{g}\) the set of \(u \in \mathfrak{gl}(V)\) such that u.x = 0 for all \(x \in \mathbf{T}(V)^{\mathfrak{g}}\) . We are going to show that \(\overline{g} = g\) .

\begin{enumerate}
\def\labelenumi{\alph{enumi})}
\item
  Show that the representation of g on the dual V* of V is isomorphic to its representation on \(\bigwedge^{p-1}V\) , where \(p = \dim V\) (use the fact that g is contained in \(\mathfrak{sl}(V)\) ). Deduce that every element of \(\mathbf{T}_{n,m} = \mathbf{T}^{n}(\mathrm{V}) \otimes \mathbf{T}^{m}(\mathrm{V}^{*})\) invariant under g is invariant under \(\overline{g}\) , and that \(\overline{g}\) is algebraic (Exerc. 5).
\item
  Let W be a vector subspace of one of the \(T_{n,m}\) . Assume that W is stable under g. Show that W is stable under \(\overline{g}\) (if \(e_{1},\ldots,e_{r}\) is a basis of W, remark that \(e_{1}\wedge\cdots\wedge e_{r}\) is invariant under g, and hence under \(\overline{g}\) ).
\end{enumerate}

Deduce that g is an ideal of \(\overline{g}\) , and that \(\overline{g}/g\) is commutative (cf.~proof of Prop. 4). We have \(\overline{g} = g \times c\) , where c is the centre of \(\overline{g}\) .

\begin{enumerate}
\def\labelenumi{\alph{enumi})}
\setcounter{enumi}{2}
\item
  Let R be the associative subalgebra of \(\mathfrak{gl}(V)\) generated by 1 and g. Show that c is contained in the centre of R (remark that \(\overline{g}\) is contained in the bicommutant of R, which is equal to R). Deduce that the elements of c are semi-simple.
\item
  Let \(x \in c\) . Show that the replicas of x belong to c (Chap. I, §5, Exerc. 14). Show that \(\operatorname{Tr}(sx) = 0\) for all \(s \in g\) ; deduce that \(\operatorname{Tr}(sx) = 0\) for all \(s \in \overline{g}\) , and hence that x is nilpotent (loc. cit).
\item
  Show that \(\mathfrak{c} = 0\) and \(\overline{\mathfrak{g}} = \mathfrak{g}\) by combining \(c\) and \(d\) ).
\end{enumerate}

\begin{enumerate}
\def\labelenumi{\arabic{enumi})}
\setcounter{enumi}{6}
\tightlist
\item
  Let \(\mathfrak{g}\) be a Lie subalgebra of \(\mathfrak{gl}(\mathrm{V})\) . Let \(m, n\) be two integers \(\geq 0\) , W and \(\mathrm{W}'\) two vector subspaces of \(\mathbf{T}^m(\mathrm{V}) \otimes \mathbf{T}^n(\mathrm{V}^*)\) where \(\mathrm{V}^*\) is the dual of V. Assume that \(\mathrm{W}' \subset \mathrm{W}\) and that W and \(\mathrm{W}'\) are stable under the natural representation of \(\mathfrak{g}\) on \(\mathbf{T}^m(\mathrm{V}) \otimes \mathbf{T}^n(\mathrm{V}^*)\) . Show that W and \(\mathrm{W}'\) are then stable under \(e(\mathfrak{g})\) . If \(\pi\) denotes the representation of \(e(\mathfrak{g})\) on \(\mathrm{W/W}'\) thus obtained, show that \(\pi e(\mathfrak{g})\) is the decomposable envelope of \(\pi(\mathfrak{g})\) (use Th. 1).
\end{enumerate}

Deduce that \(\operatorname{ad} e(\mathfrak{g})\) is the decomposable envelope of \(\operatorname{ad} \mathfrak{g}\) in \(\mathfrak{gl}(\mathfrak{g})\) .

\begin{enumerate}
\def\labelenumi{\arabic{enumi})}
\setcounter{enumi}{7}
\tightlist
\item
  Let g be a Lie subalgebra of \(\mathfrak{gl}(V)\) and h a Cartan subalgebra of g.
\end{enumerate}

\begin{enumerate}
\def\labelenumi{\alph{enumi})}
\tightlist
\item
  Show that \(e(\mathfrak{g}) = e(\mathfrak{h}) + \mathcal{D}\mathfrak{g} = e(\mathfrak{h}) + \mathfrak{g}\) .
\end{enumerate}

(Remark that \(e(\mathfrak{h}) + \mathcal{D}\mathfrak{g}\) is decomposable (Cor. 1 of Th. 1), contains \(\mathfrak{g} = \mathfrak{h} + \mathcal{D}\mathfrak{g}\) , and is contained in \(e(\mathfrak{g})\) ; thus, it is \(e(\mathfrak{g})\) .)

\begin{enumerate}
\def\labelenumi{\alph{enumi})}
\setcounter{enumi}{1}
\item
  We have \(e(\mathfrak{h}) \cap \mathfrak{g} = \mathfrak{h}\) (remark that \(e(\mathfrak{h}) \cap \mathfrak{g}\) is nilpotent).
\item
  Let \(x\) be an element of the normalizer of \(e(\mathfrak{h})\) in \(e(\mathfrak{g})\) . Show that \(x \in e(\mathfrak{h})\) . (Write \(x = y + z\) with \(y \in e(\mathfrak{h}), z \in \mathfrak{g}\) , cf.~a); remark that \([z, \mathfrak{h}] \subset e(\mathfrak{h}) \cap \mathfrak{g} = \mathfrak{h}\) , hence \(z \in \mathfrak{h}\) .)
\end{enumerate}

\(d\) ) Show that \(e(\mathfrak{h})\) is a Cartan subalgebra of \(e(\mathfrak{g})\) .

\begin{enumerate}
\def\labelenumi{\arabic{enumi})}
\setcounter{enumi}{8}
\tightlist
\item
  Let \(\mathfrak{g}\) be a Lie subalgebra of \(\mathfrak{gl}(\mathrm{V})\) . Show that conditions (i), (ii), (iii) of Exerc. 16 of §3 are equivalent to:
\end{enumerate}

\begin{enumerate}
\def\labelenumi{(\alph{enumi})}
\setcounter{enumi}{21}
\tightlist
\item
  \(\mathfrak{g}\) is decomposable and has the same rank as \(\mathfrak{g} / \mathfrak{n}_{\mathrm{V}}(\mathfrak{g})\) .
\end{enumerate}

(If \(\mathfrak{h}\) is a Cartan subalgebra of \(\mathfrak{g}\) , the condition `` \(\mathfrak{g}\) has the same rank as \(\mathfrak{g}/\mathfrak{n}_{\mathrm{V}}(\mathfrak{g})\) '' is equivalent to saying that \(\mathfrak{h} \cap \mathfrak{n}_{\mathrm{V}}(\mathfrak{g}) = 0\) , i.e.~that \(\mathfrak{h}\) has no nilpotent element \(\neq 0\) (cf.~§2, Exerc. 14). Deduce the equivalence of (i) and (v).)

\begin{enumerate}
\def\labelenumi{\arabic{enumi})}
\setcounter{enumi}{9}
\tightlist
\item
  Let \(k'\) be an extension of \(k\) and \(\mathfrak{g}'\) a \(k'\) -Lie subalgebra of
\end{enumerate}

\[
\mathfrak {g l} (\mathrm{V} \otimes_ {k} k ^ {\prime}) = \mathfrak {g l} (\mathrm{V}) \otimes_ {k} k ^ {\prime}.
\]

\begin{enumerate}
\def\labelenumi{\alph{enumi})}
\item
  Show that there exists a smallest Lie subalgebra g of \(\mathfrak{gl}(V)\) such that \(g\otimes_{k}k'\) contains \(g'\) .
\item
  Assume that \(k'\) is algebraically closed and denote by G the group of k-automorphisms of \(k'\) ; this group operates in a natural way on \(V \otimes_{k} k'\) . Show that \(g \otimes_{k} k'\) is the Lie subalgebra generated by the conjugates of \(g'\) under G (use the fact that k is the field of G-invariants in \(k'\) ).
\item
  Show that \(\mathfrak{g}\) is decomposable if \(\mathfrak{g}'\) is decomposable. (Reduce to the case where \(k'\) is algebraically closed, and use \(b\) ) as well as Cor. 1 and 3 of Th. 1.)
\end{enumerate}

¶ 11) Exceptionally, we assume in this exercise that k is a perfect field of characteristic p \textgreater{} 0.

Let g be a Lie p-algebra (Chap. I, §1, Exerc. 20). If \(x \in g\) , denote by \(\langle x \rangle\) the smallest Lie p-subalgebra containing x. It is commutative and generated as a k-vector space by the \(x^{p^{i}}\) where \(i = 0, 1, \ldots\) . Then x is said to be nilpotent (resp. semi-simple) if the p-map of \(\langle x \rangle\) is nilpotent (resp. bijective).

\begin{enumerate}
\def\labelenumi{\alph{enumi})}
\item
  Show that x can be decomposed uniquely in the form \(x = s + n\) , with \(s, n \in \langle x \rangle\) , s semi-simple and n nilpotent (apply Exerc. 23 of Chap. I, §1). If f is a p-homomorphism from g to \(\mathfrak{gl}(\mathrm{V})\) , \(f(s)\) and \(f(n)\) are the semi-simple and nilpotent components of the endomorphism \(f(x)\) ; this applies in particular to f = ad.
\item
  A subalgebra of g is said to be decomposable if it contains the semi-simple and nilpotent components of its elements. Show that, if b and c are vector subspaces of g such that \(b \subset c\) , the set of \(x \in g\) such that \([x, c] \subset b\) is decomposable (same proof as Prop. 3); in particular, every Cartan subalgebra of g is decomposable.
\item
  Let t be a commutative subalgebra of g consisting of semi-simple elements, and maximal with this property. Let h be the commutant of t in g. Let \(x \in h\) and let \(x = s + n\) be its canonical decomposition; since h is decomposable (cf.~b)), \(s, n \in h\) . Show that the subalgebra generated by t and s is commutative and consists of semi-simple elements, hence coincides with t. Deduce that \(ad_{h}x = ad_{h}n\) is nilpotent, and hence that h is nilpotent. Since \(\mathfrak{h} = \mathfrak{g}^{0}(\mathfrak{h})\) , h is a Cartan subalgebra of g (§2, Prop. 4).
\end{enumerate}

In particular, every Lie p-algebra over a finite field has a Cartan subalgebra. \(^{9}\)

